\section{Study}
\label{sec:study}

\paragraph{Tasks and models.}
Both parts below use those \StudyHeldTasks{} held-out position tasks (\StudyAlgorithms{} rules,
two variants each) on synthetic panels with \StudySeeds{} seeds, and \StudyModels{} open-weight
instruction models from \StudyFamilies{} families, 2B to 32B parameters, in BF16 at temperature
0.1. Every refusal, context stop, truncated reply and failed program stays in the denominator.
Variants and seeds are not independent of the three rules: with three clusters a sign-flip test
cannot go below $p=0.25$, so we report effects per model and pooled, without significance
claims.

\paragraph{The library on held-out tasks.}
Five arms (Figure~\ref{fig:matrix}) share the system prompt, the task text and \texttt{Check}
feedback, with at most \FixAttempts{} programs, \FixCalls{} model calls and \FixTokens{}
generated tokens per episode (\FixOutputCap{} per reply). F0 has no library. F1 adds the
position interface to the system prompt (\FixInterfaceChars{} characters, the
\FixInterfaceAlgorithms{} algorithms \texttt{positions()} accepts) and allows
\texttt{import fin\_skills}. F2 adds to the task \emph{The fin\_skills library implements this
rule as \texttt{positions("<id>", ...)}} with the id from us; one such call with the task's
parameters solves every task, so F2 shows what finding the call can give. F3 is F1 with each
algorithm's rule stated in one line. F4 first calls \texttt{route()} on the task text (reply
capped at \FixRouteCap{} tokens, charged to the episode) and continues as F2 with the id the
model named, or as F1. Every arm is judged with the \FixVisibleChecks{} checks; each has
\FixEpisodesPerArm{} episodes per model. F1$-$F0, F2$-$F0 and F2$-$F1 were declared before the
first run; F3, F4 and the split at \FixSizeCut{} parameters were added after seeing results and
are exploratory.

\paragraph{Feedback and self-improvement without the library.}
A study fixed before any test run, on the suite as it stood on 8~November 2026, ran the three
controlled arms of \S\ref{sec:methods} on the same tasks, seeds and models (\StudyOutputCap{}
generated tokens per call; \StudyPEpisodes{} episodes per model with and without the library)
and the \NMethods{} methods on \StudyCoreModels{} core models (\StudySEpisodes{} episodes,
\StudySFits{} fits; their first prompt includes the rule cards that apply). It declared
\texttt{Check} against \texttt{Repeat} and \texttt{Explain} against \texttt{Check}. We report its
runs without the library, judged with the first three checks; its library conditions used an
earlier interface and are released.
